\heading{Statistical fairness: two GPs, two different jobs}
The background GP estimates the expected counts used by the likelihood. The
significance-field GP approximates the joint null distribution of the resulting
scan. They do not need to be statistically independent: the second must reproduce
the response of the first, including its sideband training and profiling. Calling
both objects GPs neither establishes nor invalidates the significance calculation.
The substantive question is whether the generating background and the complete
analysis procedure represent the null being tested.

In this application let $r_j(n)$ be the signed profile-likelihood root at mass
$m_j$, and let $B$ denote a fixed generating spectrum. The saved response is
\begin{equation}
 r_j(n)\simeq a_j+\sum_b D_{bj}\xi_b,\qquad
 a_j=r_j(B),\quad
 D_{bj}=\sqrt{B_b}\left.\frac{\partial r_j}{\partial n_b}\right|_B,
 \quad \xi_b\sim N(0,1).
\end{equation}
Thus $s_j^2=\sum_b D_{bj}^2$ and
$R_{ij}=\sum_bD_{bi}D_{bj}/(s_i s_j)$ describe the scale and correlation of the
\emph{fitted response}. The derivatives include the response of the GP prediction
and its count-dependent covariance to fluctuations in the training bins. This is
more information than the overlap of the narrow signal templates alone.

Ananiev--Read model a signed local significance field by a zero-mean,
unit-variance GP, estimate its covariance from background response data, and
sample its maximum to estimate the trials factor. Their comparison with an
upcrossing bound concerns the same null field; it does not promise a smaller
penalty than an arbitrary resolution-based \v{S}id\'ak count. Their finite-grid
discussion also warns that coarse sampling can miss crossings and maxima
\cite{AnanievRead}. The additional $a_j,s_j$ transformation here is an HPS
reference-response construction, not a new requirement imposed by that paper.

The nominal source $B$ was fitted to the observed spectrum and then held fixed.
Estimating a background from data is familiar in a parametric bootstrap and is
not by itself a statistical error. The limitation is that this calculation has
not calibrated the outer source-fitting and method-selection procedure, or
established robustness to an absorbed signal or an inadequate background family.
Independent Poisson draws validate the field approximation \emph{conditional on
that fixed source}; they do not establish independence of the source from the
observed data. Favorable stress response must therefore be one diagnostic among
coverage, injected-signal recovery, background adequacy and held-out behavior.

\heading{Keeping the raw ordering is a separate, explicit option}
There is no obligation to replace the observed likelihood scan by
$z_j=(r_j-a_j)/s_j$ to perform a look-elsewhere calculation. Two validly specified
\emph{conditional} questions are
\begin{align}
 T_r&=\max_j\max(0,r_j), &p_{\rm global}^{(r)}&=P_B(T_r^*\ge T_r^{\rm obs}),\\
 T_z&=\max_{j:r_j>0}\frac{r_j-a_j}{s_j},
 &p_{\rm global}^{(z)}&=P_B(T_z^*\ge T_z^{\rm obs}).
\end{align}
Mass-dependent centering and scaling change the ordering, so these are different
tests of the same events. Neither test is to be chosen retrospectively because
it yields the smaller probability. Retaining $T_r$ preserves the original
peak ranking; comparing it to its simulated null still calibrates its
probability, without subtracting an offset from the observed fit.

\begin{table}[htbp]\centering\small
\begin{tabular}{lrr}
\hline
2016, 39--180 MeV & Raw maximum $T_r$ & Reference maximum $T_z$\\
\hline
Selected mass [MeV] &90.5&91.5\\
Raw root $r$ &3.452500&3.333160\\
Deterministic response $a$ &0.712849&0.496739\\
Response scale $s$ &0.982030&0.982836\\
Reference coordinate $(r-a)/s$ &2.789784&2.885955\\
Fixed-$B$ Gaussian local tail &0.002637&0.001951\\
Gaussian global exceedances / 200,000 &13,559&20,975\\
Gaussian global probability &0.067795&0.104875\\
Poisson global exceedances / 256 &18&36\\
Poisson global probability &0.0703&0.1406\\
Poisson global 95\% interval &[0.0422, 0.1088]&[0.1005, 0.1893]\\
\hline
\end{tabular}
\caption{New comparison using the same frozen v5.8.2 fields and the same 256
complete Poisson scans. No likelihood fit or background source was changed.
Intervals are Clopper--Pearson; Gaussian probabilities are conditional model
approximations. Counts, $k/N$ and add-one estimates are recorded separately in
the numerical ledger.}
\label{tab:raw-reference-statistics}
\end{table}

At 90.5 MeV, $\overline\Phi(3.452500)=0.0002777$ is the conventional standard-normal
mapping of the raw root, whereas the fixed-source response gives
$\overline\Phi[(3.452500-0.712849)/0.982030]=0.002637$.
Dividing the conditional raw-global probability by the first of these mixes a
marginal-calibration change with the look-elsewhere factor. The 256 Poisson toys
have zero local exceedances at this mass: they cannot resolve a local tail of
order $10^{-3}$, and do not prove that either precise local approximation is
correct. Passing previous statistical tests supports the procedure on the nulls
and observables those tests covered; it does not make a new source-conditioned
centering mandatory, or guarantee every marginal tail under every new source.

\heading{An apples-to-apples upcrossing comparison}
For the same standardized Gaussian grid, define
$N_u=\sum_{j=1}^{J-1}{\bf1}(Z_j\le u,Z_{j+1}>u)$. A scan exceeds $u$ only if its
first node already exceeds $u$, or there is an upcrossing. Hence
\begin{equation}
 P(\max_j Z_j>u)\le
 \overline\Phi(u)+E[N_u],\qquad
 E[N_u]=\sum_{j=1}^{J-1}
 2\,T\!\left(u,\sqrt{\frac{1-R_{j,j+1}}{1+R_{j,j+1}}}\right),
 \label{eq:finite-grid-upcrossing}
\end{equation}
where $T$ is Owen's $T$ function. The right-hand side is clipped at one. This
identity for each adjacent crossing probability gives a bound for the
\emph{actual saved Gaussian grid}; it does not assume a differentiable continuum
kernel. All thresholds shown below exceed the positive-fit gate at every mass,
so that the gate makes no change to these comparisons.

At the 2016 threshold $u=2.885955$, the saved correlated-field tail is
$0.104875$ and the same-grid upcrossing bound is $0.120378$.
The other datasets show the same ordering: $0.02345<0.02527$ (2015),
$0.33352<0.42689$ (2021), and $0.22953<0.27340$ (combined).
This is the relevant less-conservative comparison.
By contrast, the 2016 resolution-count \v{S}id\'ak value $0.02953$ and the smooth
signal-template upcrossing estimate $0.06521$ describe simplified correlation
models. The independent 283-node comparison gives $0.42461$.
The full fitted-response field lies between these illustrative answers; the
smaller resolution-only answers do not establish overconservatism of that field.

\fig[1]{stats_same_field_lee}{New 2016 comparison with the source and
test held fixed. Left: the sampled Gaussian maximum is less conservative than
the analytic upcrossing bound for the same finite grid. Right: counting
representative points misses searches within the intervening regions; treating
the properly calculated region maxima as independent errs in the other
direction in this example. The left curve reuses 200,000 saved fields; the right
uses 500,000 fresh fields with the same covariance. Shading denotes binomial
95\% intervals for the joint maxima. These are fixed-background diagnostics.}

\heading{What a 10\% correlation criterion can and cannot do}
For $N$ independent, identically calibrated point tests,
$p_{\rm global}=1-(1-p_{\rm local})^N$. A tolerance $|\rho|<0.1$ is not
independence, and selecting weakly correlated points does not preserve the
maximum of the masses omitted between them. The saved 2016 field makes both
issues concrete. A left-to-right rule that selects the next mass with
$|\rho|<0.1$ relative to the previous selected mass yields 24 points. Their
non-neighbor correlations still reach $|\rho|=0.835$.
Requiring the bound against \emph{every} selected point yields six points when
starting from the low-mass end and seven when starting from the high-mass end.
The count depends on the selection rule. Every adjacent native-grid correlation
is at least $0.889$, so no split into two contiguous sets makes all correlations
across the boundary smaller than 0.1.

A defensible regional version would first compute
\begin{equation}
 p_k(u)=P\!\left(\max_{m\in\mathcal R_k} Z(m)>u\right),\qquad
 p_{\rm independent\ regions}(u)=1-\prod_k[1-p_k(u)],
\end{equation}
and then test independence of the \emph{regional maxima}. The maximum within
each region must already include its internal look-elsewhere effect.
Using midpoint blocks around the 24 representative masses, the new conditional
simulation gives a full-grid tail $53{,}201/500{,}000=0.10640$
(95\% interval $[0.10555,0.10726]$). Multiplying the regional non-exceedance
probabilities gives $0.14341$; pretending each region is one normal point gives
$0.04579$. The six-point all-pairs rule gives $0.01165$.
None is a substitute for the already available joint-field maximum.

\fig[.94]{stats_correlation_selection}{The proposed threshold made
explicit. Left: low correlation with the immediately preceding selected point
does not prevent strong correlations with more distant selected points.
Right: 24 representative points do not capture the full observed scan or its
maximum. Vertical lines delimit one illustrative contiguous partition, not
independently established signal regions.}

The tail-equivalent count
$N_{\rm eff}(u)=\log[1-p_{\rm global}(u)]/\log[1-\overline\Phi(u)]$
is a useful way to quote the answer from the full field; it is generally
threshold dependent. For the saved 2016 peak it is $56.73$, whereas the spectral
participation rank is $9.16$ and the grid contains 283 nodes. These measure
different things. For a smooth stationary field, the high-threshold crossing
term scales as $\exp(-u^2/2)$ while the local tail scales approximately as
$\exp(-u^2/2)/u$, already explaining why a single constant effective count need
not describe all thresholds. A discovery-threshold penalty also does not, by
itself, change a fixed-mass exclusion limit; a new discovery reach requires
signal-injection power for the specified global test.

The initial 100,000-draw block check had an unusually high tail of $0.10910$.
The covariance reconstruction and positive-fit gate were verified, and the
same seed was extended to 500,000 draws. The result above differs from the saved
200,000-draw estimate by about $1.9$ combined Monte Carlo standard errors; their
95\% intervals overlap. This sampling check is recorded explicitly rather than
being interpreted as a change to the physical model or a calibrated source
uncertainty.
